\section*{Part VII: The Bayes Institute Beachhead Decision}

\subsection*{85. Begin with What Is Actually True}

Bayes Institute's durable purpose is to create the conditions in which human possibility becomes human capability. Its campaign, ``BECOME MORE.,'' expresses a long-range aspiration. Its intended field is rigorous STEM education with strong links to technology, industry, and eventually research. It is not intended to remain forever a short-course seller, nor to acquire university status by declaring itself one.

The founder has a technical background spanning software engineering, applied AI, statistics, teaching, and curriculum design, plus some experience talking with colleges and industry organizations. The founder also intends to retain a full-time professional role and initially operate with a lean cash budget and constrained weekly hours. These are meaningful starting assets and constraints.

What Bayes does \emph{not} yet possess should be equally explicit: a validated repeatable cohort-acquisition engine, a longitudinal Bayes alumni record, dependable independent faculty capacity, confirmed large institutional contracts, a mature research operation, a degree-granting authorization, or a physical campus. Other personal experience and earlier projects should not be presented as Bayes institutional outcomes.

The beachhead decision must fit both sides of that ledger. Founder expertise can justify a difficult technical demonstration. It cannot automatically support an untested claim of employment transformation or global institutional recognition.

\subsection*{86. Preserve Ambition Without a Premature Product List}

The future institution could plausibly contain academic programs, research groups, applied laboratories, college and industry partnerships, enterprise capability building, incubation, and a physical campus. But each is a different operational system with different capital, governance, and legitimacy requirements.

For the first market, consider 3 outcomes separately:

\begin{itemize}
    \item \textbf{Cash now:} Can the offer produce real, collectible contribution without unacceptable risk?
    \item \textbf{Capabilities next:} Does it build faculty standards, instructional assets, real technical projects, employer-relevant assessment, relationships, or repeatable infrastructure?
    \item \textbf{Institutional integrity:} Does it reinforce the standards and claims that a future serious institute should be proud to inherit?
\end{itemize}

The first does not guarantee the second. The second does not excuse the absence of the first. The third can constrain both, especially when a short-term revenue opportunity depends on misleading career or accreditation claims.

\subsection*{87. There Are At Least 3 Genuinely Distinct Initial Customer Systems}

Three separate customer groups present themselves:

\begin{itemize}
    \item \textbf{D2C professional learners:} Purchase using personal funds and scarce time. Their buying trigger may be a career inflection or a live engineering responsibility. They compare free resources, peers, mentors, bootcamps, and self-study. Evidence of rigorous feedback may matter more than a brochure.
    \item \textbf{Colleges and universities:} May purchase applied electives, faculty-enabled curriculum, project labs, or complementary career services. The user is often the student; the economic buyer and decision authority may be the college, department, management, or some combination. Semester calendars, approvals, timetable fit, faculty coordination, and payment terms strongly influence the sale.
    \item \textbf{Enterprises:} May purchase team AI training, diagnostics, workflow improvement, or applied technical advisory services. The user is the employee; the budget owner may be an engineering head, HR/L\&D team, or business leader. The purchasing decision depends on a current workflow, business consequences, procurement, security, and evidence of delivery.
\end{itemize}

A founder can eventually participate in all 3. But they cannot sensibly use the same market size, conversion funnel, product promise, or operating process for all 3.

\subsection*{88. Candidate A: Working Engineers Facing Consequential AI Tasks}

Imagine a 24 to 30-year-old developer with real software experience who has been assigned to lead or materially contribute to a retrieval, agentic workflow, model-evaluation system, or other AI-enabled product. They can write code and use modern tools but lack confidence in evaluation design, reliability, security boundaries, architecture trade-offs, and the ability to defend their decisions.

The job is not ``learn AI'' generally. It is \textbf{become able to design, evaluate, and explain an AI system responsibly under practical constraints.} The alternative may be trial and error, existing senior engineers, public documentation, community help, or a company-provided learning budget.

Potential founder advantage: deep software/AI knowledge and substantial teaching experience. Potential wedge: inspectable technical review against a demanding rubric, culminating in independently defended work. Potential entry product: a short paid technical studio or diagnostic with a complete, bounded outcome.

Uncertainties are substantial. The developer may regard coaching as unnecessary, may lack time, may need employer funding, or may prefer one-off feedback to a sustained course. The founder must discover whether this is a customer group with repeatable purchasing behavior, not simply a technically attractive target.

\subsection*{89. Narrow Candidate A Further by Trigger, Not Aspiration}

Compare 4 descriptions:

\begin{enumerate}
    \item ``People aged 21 to 30 interested in AI.''
    \item ``Working software engineers who want to learn AI.''
    \item ``Backend engineers who are building AI features at work.''
    \item ``Backend engineers responsible for an AI system that must pass a release or architecture review within roughly 1 quarter, who lack adequate internal review capacity.''
\end{enumerate}

Each step adds specificity. The fourth description may be easiest to connect to urgency and an observable technical job, but it may be too narrow or too difficult to recruit. It might also be that employed engineers cannot share relevant project details. The right definition is empirical: narrow enough for a clear buying and delivery mechanism, broad enough to support repeated opportunities.

The founder should interview those who recently faced the trigger, including people who solved it without training. Their alternatives provide a better price and product reference than competing course syllabi alone.

\subsection*{90. The First Paid Product Should Test the Wedge}

Suppose Bayes offers a \textbf{3-week applied evaluation studio} (hypothetical design). Eligibility requires solid programming skills. The learner receives a deliberately flawed example AI workflow, a technical rubric, scheduled expert review, and a final defense of corrected architecture and evaluation evidence. Learners work on synthetic or approved non-confidential material, not unauthorized company data.

The promise is bounded: the program will offer specific instruction, reviewed submissions, and a verifiable assessment process. It does not promise a job, promotion, production deployment, or universally transferable mastery from 3 weeks of practice.

What this experiment could test: whether such engineers find the problem urgent, will pay a real fee, undertake demanding work, value high-quality feedback, and improve on the specified task. It might also reveal review labor, reliable scheduling, and delivery cost.

What it would \emph{not} establish: the viability of a 6-month flagship, large-scale acquisition, employer endorsement, long-term career effect, or capability transfer across every AI domain.

\subsection*{91. A Competing D2C Wedge: Personal Capability Audit}

A second possible wedge is a paid, 2-session engineering capability diagnostic. Learners attempt a novel design or evaluation task and receive candid expert analysis of performance gaps and a recommended learning plan.

This is smaller, easier to buy, and potentially an effective discovery device. It may also be \textbf{too much of a consulting service}: reviewers must spend scarce time, learners may not want repeated purchases, and an audit can reveal capability without developing it.

If it is used as a route to a longer program, evaluate both products independently. The diagnostic should create real value at its disclosed price even if the learner never enrolls. The long program must earn separate evidence of willingness to pay and deliverability.
